\documentclass[12pt]{article}
\usepackage[utf8]{inputenc}
\usepackage[T1]{fontenc}
\usepackage{hyperref}
\title{Southern Kurdish Language Resources for K-12 Curriculum Development}
\author{Compiled by Claude Sonnet 5.5 (claude-sonnet-5-5) for Andr\'as Kornai}
\date{}
\begin{document}
\maketitle

\section{Introduction}
Southern Kurdish (ISO 639-3 sdh; Glottocode sout2640) is spoken by about three million people in western Iran (notably Kermanshah and Ilam provinces) and in Iraq. Its varieties include Kalhori, Feyli, Garrusi and Kordali. Laki is sometimes counted as a separate language and sometimes as part of the group. No school curriculum in Southern Kurdish was found. The resources located are research-oriented.

\section{Primary Educational Resources}
No primers or graded textbooks for Southern Kurdish were found. The only Kurdish school-textbook resource that turned up was a corpus of 31 K-12 textbooks in Sorani, which is a different variety. See the corpus paper below.

\section{Linguistic and Reference Materials}
\begin{itemize}
\item Wikipedia article on Southern Kurdish, with dialect classification and references: \url{https://en.wikipedia.org/wiki/Southern_Kurdish}
\item Wikipedia article on Laki: \url{https://en.wikipedia.org/wiki/Laki_language}
\item eLex 2021 paper describing a machine-readable Southern Kurdish lexicon built from dictionaries (reported to be released under CC BY-SA 4.0): \url{https://elex.link/elex2021/wp-content/uploads/2021/08/eLex_2021_30_pp479-488.pdf}
\item Per the Wikipedia article, Sara Belelli's 2021 monograph \emph{The Laki Variety of Harsin: Grammar, Texts, Lexicon} is a reference grammar of one Laki variety (no direct URL found).
\end{itemize}

\section{Religious and Community Resources}
No scripture translations or community organisations for Southern Kurdish were found in search results.

\section{Digital and Audio Resources}
\begin{itemize}
\item Corpus creation for Southern Kurdish and Laki (arXiv, 2023): \url{https://arxiv.org/pdf/2304.01319}
\item A fine-grained Kurdish corpus paper (the Sorani KTC textbook corpus, not Southern Kurdish): \url{https://arxiv.org/pdf/1909.11467}
\end{itemize}

\section{Curriculum Development Context}
Southern Kurdish is not known to be taught in schools. Persian is the language of instruction in Iran and Arabic or Sorani in Iraq. Orthographic standardisation, graded readers and teacher materials do not appear to exist. The corpora and lexicon above could support the development of such materials. Searches in Persian and Kurdish for Kalhori and Laki teaching resources found nothing relevant.
\end{document}